\documentclass{beamer}
\usetheme{Madrid}
\usecolortheme{default}

\usepackage[utf8]{inputenc}
\usepackage[T1]{fontenc}
\usepackage[french]{babel}
\usepackage{xcolor}
\usepackage{tikz}
\usepackage{fontawesome}

% Custom colors
\definecolor{hydroblue}{RGB}{0,102,204}
\definecolor{crimsonred}{RGB}{220,20,60}
\definecolor{forensicgreen}{RGB}{34,139,34}

\title[Mobile Forensics]{Analyse Forensique Samsung Galaxy S10}
\subtitle{DFRWS Challenge 2021 - Étape 3}
\author{Justin Bossard \and Mattéo Gouhier \and Paul Mathé \and Samuel Plet}
\institute{UQAC - Atelier Pratique en Cybersécurité}
\date{Octobre 2025}

\begin{document}

% ============= SLIDE 1: TITRE =============
\begin{frame}
\titlepage
\end{frame}

% ============= ACTE 1: L'ENQUÊTE COMMENCE (5 min) =============
% Justin (1.5 min) → Samuel (2 min) → Paul (1.5 min)

% SLIDE 2: Contexte mystérieux (Justin - 1.5 min)
\begin{frame}{\faSearch \, Le Mystère de Lausanne}
\begin{center}
\Large{\textcolor{crimsonred}{Lausanne, 9 avril 2021}}
\end{center}

\begin{block}{Un crime découvert}
\begin{itemize}
    \item \faMapMarker \, ATM compromis - Dispositif de skimming trouvé
    \item \faCalendar \, 11 jours plus tard : Arrestation mystérieuse à Genève
    \item \faQuestion \, \textbf{Qui est le suspect ?}
\end{itemize}
\end{block}

\begin{alertblock}{Notre Mission}
\textbf{DFRWS Challenge 2021} : Analyser un Samsung Galaxy S10 pour élucider l'affaire
\end{alertblock}

\pause
\begin{center}
\Large{\textcolor{hydroblue}{Samuel va nous révéler qui se cache derrière ce crime...}}
\end{center}
\end{frame}

% SLIDE 3: Révélation identité (Samuel - 2 min)
\begin{frame}{\faUser \, Révélation de l'Identité}
\begin{columns}
\begin{column}{0.5\textwidth}
\begin{block}{Autopsy révèle...}
\begin{itemize}
    \item \faInstagram \, Profil Instagram découvert
    \item \faUser \, \textbf{Werner von Strucker}
    \item \faAt \, \texttt{therealredskull41@gmail.com}
\end{itemize}
\end{block}

\begin{alertblock}{Corrélation}
Photos Instagram $\leftrightarrow$ Photos téléphone
\end{alertblock}
\end{column}

\begin{column}{0.5\textwidth}
\begin{center}
% Placeholder pour photo du suspect
\fbox{\begin{minipage}{3cm}
\centering
\textcolor{gray}{Photo\\Werner\\von Strucker}
\end{minipage}}
\end{center}
\end{column}
\end{columns}

\pause
\begin{center}
\Large{\textcolor{crimsonred}{HYDRA... Paul, ce nom te dit quelque chose ?}}
\end{center}
\end{frame}

% SLIDE 4: Lien HYDRA (Paul - 1.5 min)
\begin{frame}{\faExclamationTriangle \, Le Lien HYDRA}
\begin{block}{Paul confirme}
\begin{itemize}
    \item \faExclamationTriangle \, \textbf{HYDRA} : Organisation criminelle (Marvel)
    \item \faWarning \, "Red Skull" dans l'email
    \item \faShield \, Mais ici, ce n'est pas de la fiction...
\end{itemize}
\end{block}

\begin{center}
\tikz{
    \node[circle, fill=crimsonred, text=white, font=\Large] {H};
    \node[right=0.1cm, circle, fill=crimsonred, text=white, font=\Large] {Y};
    \node[right=0.2cm, circle, fill=crimsonred, text=white, font=\Large] {D};
    \node[right=0.3cm, circle, fill=crimsonred, text=white, font=\Large] {R};
    \node[right=0.4cm, circle, fill=crimsonred, text=white, font=\Large] {A};
}
\end{center}

\pause
\begin{center}
\Large{\textcolor{hydroblue}{Mattéo, que nous révèlent ses déplacements ?}}
\end{center}
\end{frame}

% ============= ACTE 2: LA TRAQUE (7 min) =============
% Samuel (2.5 min) → Mattéo (2.5 min) → Mattéo suite (1.5 min) → Paul (0.5 min)

% SLIDE 5: Géolocalisation (Samuel - 2.5 min)
\begin{frame}{\faMapO \, La Traque Géographique}
\begin{columns}
\begin{column}{0.6\textwidth}
\begin{block}{16 Photos Géolocalisées}
\begin{itemize}
    \item \faMapMarker \, Région lausannoise
    \item \faHome \, Cache Google Maps :
    \begin{itemize}
        \item Bussigny
        \item Prilly  
        \item Renens
    \end{itemize}
    \item \faTrain \, 15 avril : Voyage à Genève
\end{itemize}
\end{block}
\end{column}

\begin{column}{0.4\textwidth}
\begin{center}
% Placeholder pour carte
\fbox{\begin{minipage}{3cm}
\centering
\textcolor{gray}{Carte\\Suisse\\Déplacements}
\end{minipage}}
\end{center}
\end{column}
\end{columns}

\pause
\begin{alertblock}{Voyage Suspect}
\faCalendar \, 15 avril : Gare → ICC Genève
\end{alertblock}

\pause
\begin{center}
\Large{\textcolor{crimsonred}{Mattéo a découvert des messages secrets...}}
\end{center}
\end{frame}

% SLIDE 6: Messages secrets (Mattéo - 2 min)
\begin{frame}{\faLock \, Messages Cryptés Découverts}
\begin{block}{OmniNotes Suspect}
\begin{itemize}
    \item \faFileTextO \, Application de notes avec contenu chiffré
    \item \faKey \, Préfixe mystérieux récurrent : \textcolor{crimsonred}{\texttt{"HHY"}}
    \item \faCogs \, Nécessité de reverse engineering
\end{itemize}
\end{block}

\begin{center}
\fbox{\begin{minipage}{8cm}
\ttfamily\small
HHYiVGhhdCBpcyBub3QgZGVhZCB3aGljaCBjYW4gZXRlcm5hbA==\\
HHYQWxsIGhhaWwgdGhlIG1pZ2h0eSBIeWRyYSEgUmVkIFNrdWxs\\
HHYMeXN0ZXJpb3VzIG1lc3NhZ2UgZnJvbSB0aGUgZGFyayBzaWRl
\end{minipage}}
\end{center}

\pause
\begin{center}
\Large{\textcolor{hydroblue}{Devinez ce que signifie HHY ?}}
\end{center}
\end{frame}

% SLIDE 7: Déchiffrement (Mattéo suite - 1.5 min)
\begin{frame}{\faUnlock \, Algorithme Révélé}
\begin{block}{Reverse Engineering Réussi}
\begin{itemize}
    \item \faCog \, \textbf{AES/ECB} + Base64 + préfixe "HHY"
    \item \faKey \, Clé hybride découverte : \texttt{rb5CuefkDLyb9T} + date
    \item \faCalendar \, Format date : YYYY-MM-DD
    \item \faUnlock \, \textbf{42 messages décodés}
\end{itemize}
\end{block}

\pause
\begin{alertblock}{Révélation Choc}
\centering
\Large{\textcolor{crimsonred}{"Heil Hydra [HYDRA]!"}}
\end{alertblock}

\pause
\begin{center}
\Large{\textcolor{forensicgreen}{Et le 15 avril ? Paul peut nous aider...}}
\end{center}
\end{frame}

% SLIDE 8: Corrélation (Samuel - 1 min)
\begin{frame}{\faBullseye \, Corrélation Parfaite}
\begin{columns}
\begin{column}{0.5\textwidth}
\begin{block}{Message du 15 avril}
\textcolor{crimsonred}{\texttt{"below Turkish consulate"}}
\end{block}

\begin{block}{Photos Genève}
\begin{itemize}
    \item \faCamera \, ICC (International Conference Center)
    \item \faFlag \, Consulat turc visible
    \item \faMapMarker \, Géolocalisation exacte
\end{itemize}
\end{block}
\end{column}

\begin{column}{0.5\textwidth}
\begin{center}
% Placeholder pour photos corrélation
\fbox{\begin{minipage}{3.5cm}
\centering
\textcolor{gray}{Photos\\ICC \& Consulat\\Corrélation}
\end{minipage}}
\end{center}
\end{column}
\end{columns}

\pause
\begin{center}
\Large{\textcolor{hydroblue}{Justin, parlons méthodologie !}}
\end{center}
\end{frame}

% ============= ACTE 3: LES PREUVES S'ACCUMULENT (5 min) =============
% Justin (1.5 min) → Paul (2 min) → Paul suite (1.5 min)

% SLIDE 9: Principe de Locard (Justin - 1.5 min)
\begin{frame}{\faBalanceScale \, Principe de Locard}
\begin{block}{Méthodologie Forensique}
\textcolor{forensicgreen}{\Large{"Chaque interaction laisse une trace"}}
\end{block}

\begin{itemize}
    \item \faCheckSquareO \, \textbf{5 phases appliquées} : Identification → Collection → Examination → Analysis → Presentation
    \item \faShield \, \textbf{Intégrité préservée} : Vérification SHA2-256
    \item \faGavel \, \textbf{Recevabilité juridique} garantie
\end{itemize}

\pause
\begin{center}
\Large{\textcolor{crimsonred}{Paul, quelles autres traces avons-nous ?}}
\end{center}
\end{frame}

% SLIDE 10: Piste financière (Paul - 2 min)
\begin{frame}{\faDollar \, L'Argent du Crime}
\begin{block}{Bitcoin Cash Découvert}
\begin{itemize}
    \item \faBitcoin \, Adresse principale : \texttt{qpuv-d23y}
    \item \faRandom \, Réseau d'obfuscation complexe
    \item \faCalendar \, \textbf{18 avril 2021} : Toutes les transactions
    \item \faDollar \, Montants faibles (\~4-5\$) pour masquer l'origine
\end{itemize}
\end{block}

\begin{center}
% Placeholder pour schéma blockchain
\fbox{\begin{minipage}{6cm}
\centering
\textcolor{gray}{Schéma\\Réseau Bitcoin\\Obfuscation}
\end{minipage}}
\end{center}

\pause
\begin{center}
\Large{\textcolor{hydroblue}{Mais l'acte criminel lui-même ?}}
\end{center}
\end{frame}

% SLIDE 11: Skimming prouvé (Paul - 1.5 min)
\begin{frame}{\faVolumeUp \, L'Audio qui Condamne}
\begin{block}{Enregistrement Suspect}
\begin{itemize}
    \item \faVolumeUp \, \texttt{2021\_04\_09T1621} : Audio mystérieux
    \item \faCalendar \, 9 avril 2021 : Corrélation temporelle parfaite
    \item \faCreditCard \, Dispositif "audioskim" physique retrouvé
\end{itemize}
\end{block}

\begin{alertblock}{Preuve Physique}
\faMapMarker \, ATM Lausanne + Dispositif = Crime prouvé
\end{alertblock}

\pause
\begin{center}
\Large{\textcolor{crimsonred}{Timeline finale à 4 voix !}}
\end{center}
\end{frame}

% ============= FINAL: TIMELINE COLLABORATIVE (3 min) =============
% Paul introduit → Tous ensemble → Justin clôture

% SLIDE 12: Timeline collaborative
\begin{frame}{\faCalendar \, Reconstitution Collaborative}
\begin{center}
\textbf{\Large{Timeline du Crime}}
\end{center}

\begin{itemize}
    \item[\faUser] \textbf{Justin} : \textcolor{hydroblue}{8 avril : Recrutement HYDRA}
    \item[\faVolumeUp] \textbf{Samuel} : \textcolor{forensicgreen}{9 avril : Skimming à Lausanne}  
    \item[\faMapMarker] \textbf{Mattéo} : \textcolor{orange}{15 avril : Rendez-vous secret Genève}
    \item[\faBitcoin] \textbf{Paul} : \textcolor{crimsonred}{18 avril : Blanchiment Bitcoin Cash}
\end{itemize}

\pause
\begin{center}
\Large{\textcolor{crimsonred}{\textbf{20 avril : Arrestation - FIN DE PARTIE !}}}
\end{center}

\pause
\begin{block}{Conclusion Forensique}
\textbf{Chaîne de preuves complète} → Intégrité préservée → \textcolor{forensicgreen}{Recevabilité juridique}
\end{block}
\end{frame}

% SLIDE 13: Questions
\begin{frame}{\faQuestion \, Questions \& Discussion}
\begin{center}
\Large{\textbf{Questions ?}}
\end{center}

\begin{block}{Expertise par domaine}
\begin{itemize}
    \item \textcolor{hydroblue}{\textbf{Justin}} : Méthodologie forensique, aspects juridiques
    \item \textcolor{forensicgreen}{\textbf{Samuel}} : OSINT, blockchain, géolocalisation  
    \item \textcolor{orange}{\textbf{Mattéo}} : Cryptographie, reverse engineering
    \item \textcolor{crimsonred}{\textbf{Paul}} : Audio forensics, skimming techniques
\end{itemize}
\end{block}

\vfill
\begin{center}
\textit{Merci pour votre attention !}
\end{center}
\end{frame}

% SLIDES DE BACKUP (si temps/questions)
% Note: Format dynamique avec interventions croisées selon la répartition
\begin{frame}{Backup: Détails Techniques AES}
\begin{block}{Algorithme de Chiffrement}
\begin{itemize}
    \item \textbf{Mode} : AES-128-ECB (Electronic Codebook)
    \item \textbf{Encodage} : Base64 après chiffrement
    \item \textbf{Clé dynamique} : \texttt{rb5CuefkDLyb9T} + date du message
    \item \textbf{Préfixe} : "HHY" ajouté avant encodage Base64
\end{itemize}
\end{block}

\begin{center}
\ttfamily\small
\fbox{\begin{minipage}{8cm}
plaintext → AES-ECB(key+date) → Base64 → "HHY"+encoded
\end{minipage}}
\end{center}
\end{frame}

\begin{frame}{Backup: Réseau Bitcoin Cash}
\begin{center}
% Placeholder pour diagramme détaillé
\fbox{\begin{minipage}{10cm}
\centering
\textcolor{gray}{Diagramme détaillé\\Réseau d'obfuscation\\Bitcoin Cash}
\end{minipage}}
\end{center}
\end{frame}

\end{document}